% CAPÍTULO: CASOS MÁS FRECUENTES

\section*{Casos más frecuentes}
Si el cable del impulsor regresó en su totalidad, será indicativo de que la fuente se encuentra atrapada dentro de la manguera. Si el equipo se desarmó antes de cerciorarse de que la fuente se encontraba dentro del contenedor y no se localiza en la manguera, lo más probable es que esté en el piso.

Si la fuente se encuentra atrapada en la manguera sin estar atorada, ésta se encontrará en o cerca de la punta. En este caso la fuente se puede recuperar si se deposita en el suelo:

\begin{itemize}
    \item Para depositar la fuente en el suelo se procederá a:
    \begin{enumerate}
        \item Llevar la fuente a la punta de la manguera con la ayuda del impulsor (verifique la posición de la fuente con el detector).
        \item Ponga la manguera a nivel del piso. Si se encuentra dentro de una tubería o algún otro artefacto, sáquela a un espacio libre que le permita efectuar las maniobras. Evite en todo momento que el contenedor sufra algún daño.
        \item Coloque el blindaje semicilíndrico sobre la manguera y junto al contenedor.
        \item Jale el equipo hasta que la fuente quede debajo del blindaje.
        \item Amarrar la manguera a uno de los extremos de la caña.
        \item Corte la manguera a 50 cm del lugar donde la amarró.
        \item Saque la manguera del blindaje y levántela con ayuda de la caja y accesorios para que la fuente caiga al suelo.
    \end{enumerate}
\end{itemize}
